\input format.tex

%\usepackage{arydshln}

\begin{document}

%%%
\setlength{\arrayrulewidth}{0.5pt}
\color{gray2}
\vspace*{0mm}
\fontsize{9.3pt}{9.8pt}\selectfont
\begin{tctabularx}{tabularx={p{15.4cm}<{\centering}p{0pt}}}
\\[-6.769pt]
\cellcolor{kelly} \raisebox{1pt}{\color{black70}\xiaosihao\lantxh\bf 肠道调节方案} \\
\end{tctabularx}

{
\lantxh
\vspace*{-4.25mm}
\noindent
\begin{longtable}{|@{~~~}p{15.35cm}|}
\hline
\parbox[c]{\hsize}{\vskip7pt\xiaowuhao {\lantxh\bf 补充益生菌、益生元和膳食纤维}\\1.每天服用适量益生菌产品（如乳双歧杆菌活菌制剂等），增加肠道内益生菌的数量，抑制有害菌的生长。

2.每天可服用适量低聚糖类产品，如低聚果糖、低聚寡糖、低聚异麦芽糖，促进体内双歧杆菌、乳杆菌等益生菌的生长。

3.每天选择性服用膳食纤维，这将有益于您的肠道健康。

注：如患炎症性肠病、腹泻，不宜服用菊粉及低聚果糖，建议选择低聚半乳糖。
\vskip7pt}\\
\hline
\parbox[c]{\hsize}{\vskip7pt\xiaowuhao {\lantxh\bf 补充营养品}\\可选择性补充复合蛋白固体饮料、大豆肽蛋白粉、含有几丁聚糖的保健品，帮助您控制体重。
\vskip7pt}\\
\hline
\end{longtable}
}


\fontsize{9.3pt}{9.8pt}\selectfont
\begin{tctabularx}{tabularx={p{15.4cm}<{\centering}p{0pt}}}
\\[-6.769pt]
\cellcolor{kelly} \raisebox{1pt}{\color{black70}\xiaosihao\lantxh\bf 运动方案} \\
\end{tctabularx}

{
\lantxh
\vspace*{-4.25mm}
\noindent
\begin{longtable}{|p{1.25cm}@{}p{2cm}|p{11.732cm}|}
\hline
\parbox[c]{\hsize}{\vskip7pt\raisebox{-1.5ex}{\includegraphics[width=.9cm]{yundong.pdf}}\vskip7pt} & \parbox[c]{\hsize}{\centering\xiaowuhao\vskip7pt{运动}\vskip7pt} &
\parbox[c]{\hsize}{\vskip7pt\xiaowuhao {1.条件允许的情况下，选择爬楼梯，自行车出行等方式锻炼身体。多参加团体运动，培养运动习惯。

2.建议选用持续时间长、低强度的运动方式，如快步走、慢跑、骑自行车、游泳、乒乓球、羽毛球、网球、健身操等。配合增加肌肉量的运动，如俯卧撑、哑铃、深蹲等。

3.建议早饭或晚饭前空腹状态下进行运动，每次运动时间30分钟以上，每周5次以上。}\vskip7pt}\\
\hline
\end{longtable}
}
\end{document}
